\subsection{Reusable validation-integrity profile}
\label{sec:vip}

We formalized the audit as the seven-domain validation-integrity profile
(Table~\ref{tab:vip}). The profile is a protocol, not a numerical score. Each domain is
reported as \emph{verified}, \emph{violated}, \emph{indeterminate}, or
\emph{not applicable}, together with executable evidence. This prevents a
high aggregate score from concealing one claim-invalidating failure.

\begin{table}[H]
\centering
\scriptsize
\caption{Reusable Validation-Integrity Profile for longitudinal prediction AI.}
\label{tab:vip}
\begin{tabular}{L{1.1cm} L{2.8cm} L{4.1cm} L{4.0cm} L{3.2cm}}
\toprule
\textbf{Code} & \textbf{Domain} & \textbf{Validity condition} & \textbf{Executable audit} & \textbf{Required evidence} \\
\midrule
E & Endpoint and index time & The target, prediction time, risk set, and predictor availability are unambiguous. & Reconstruct episodes and observation spacing; reject future targets without continuous follow-up. & Endpoint contract; exclusions; sampling diagram. \\
I & Identity separation & Every correlated identity unit belongs to only one side of an evaluation split. & Count participant overlap in every fold; compare row-wise with grouped out-of-fold estimates. & Group definition; split manifest; paired metric change. \\
T & Transformation locality & Imputation, scaling, normalization, and representation statistics use training data only. & Replace label-conditioned and test-transductive transformations sequentially. & Fitted-statistic provenance; protocol ladder. \\
S & Supervised-selection locality & Every outcome-informed selection or tuning operation is repeated inside training folds. & Compare global with fold-local selection under an otherwise fixed grouped design. & Selected features by fold; paired metric change. \\
C & Calibration integrity & Calibration units are independent of fitting/test units and reflect the target prevalence. & Measure fit--calibration identity/row overlap; compare balanced with natural-prevalence calibration. & Calibration split; prevalence; Brier, log loss, intercept, slope. \\
H & History independence & Test histories do not duplicate or substantially overlap training/evaluation histories beyond the deployment scenario. & Evaluate all prespecified non-overlapping temporal offsets or forward-chaining splits. & Window map; offset results; no performance-selected offset. \\
O & Observation-process resistance & Administrative availability or post-event proxies cannot determine the label. & Prohibited-feature stress test plus refitted within-identity permutation and circular-shift controls. & Feature-availability contract; null distributions. \\
\bottomrule
\end{tabular}
\end{table}

Application follows five steps: define the intended claim and endpoint contract;
map identities, time windows, and all outcome-informed operations; construct the
safeguarded reference pipeline; vary one integrity domain at a time while
holding the model family and evaluation rows fixed where possible; and report
paired cluster-aware changes for discrimination and probability accuracy. For
metric $m$ and stress condition $d$, the audit quantity is
\begin{equation}
\Delta_{m,d}=m(\widehat{p}_{d})-m(\widehat{p}_{\mathrm{reference}}),
\end{equation}
where both predictions are aligned to the same evaluation rows and uncertainty
is resampled at the independent identity unit. The sign is interpreted in the
metric's direction; the profile does not assume that effects are additive or causal.
